%=========== ΛΥΣΗ - ΘΕΜΑ Β_89 ===========
\begin{thema}{Β}
Στο παρακάτω αναβαθμισμένο σχήμα έχει προστεθεί διδακτικός σχολιασμός (ασύμπτωτες και κρίσιμα σημεία) για τη διευκόλυνση και πληρέστερη κατανόηση της λύσης.
\begin{center}
\begin{tikzpicture}
    \begin{axis}[
        axis lines = middle, 
        axis line style = {->, >=latex}, 
        grid = major, 
        grid style = {very thin, gray!30}, 
        xlabel = {$x$}, 
        ylabel = {$y$}, 
        xtick = {-2,1,2,3,4,5}, 
        ytick = {1,2,3,4}, 
        every axis x label/.style = {at={(current axis.right of origin)}, anchor=west}, 
        every axis y label/.style = {at={(current axis.above origin)}, anchor=south}, 
        width = 10cm, 
        height = 7.5cm, 
        samples = 100, 
        xmin=-3.5, 
        xmax=6.5, 
        ymin=-0.5, 
        ymax=4.5, 
        xtick distance=1, 
        ytick distance=1
    ]
        % Αρχικές καμπύλες
        \addplot[very thick, \xrwma, domain=-3.5:1.75] {1/((x-2)^2)};
        \addplot[very thick, \xrwma, domain=2.25:6.2] {1 + 1/((x-2)^2)};
        
        % Ασύμπτωτες
        \addplot[dashed, thick, color=red] coordinates {(2,-0.5) (2,4.5)};
        \addplot[dashed, thick, color=blue] coordinates {(2.25,1) (6.5,1)};
        
        % Annotations
        % Vertical asymptote x=2
        \node[red, anchor=west] at (axis cs:2.1, 4) {\footnotesize \textbf{Κατακόρυφη Ασύμπτωτη $x=2$}};
        \draw[->, >=stealth, red, thick] (axis cs:2.05, 3.8) -- (axis cs:2.05, 4.4);
        
        % Horizontal asymptote y=1
        \node[blue, anchor=south] at (axis cs:4.5, 1.2) {\footnotesize \textbf{Οριζόντια Ασύμπτωτη $y=1$}};
        \draw[->, >=stealth, blue, thick] (axis cs:4.5, 1.2) -- (axis cs:4.5, 1.05);

        % Horizontal asymptote y=0
        \node[black, anchor=south] at (axis cs:-1.5, 0.4) {\footnotesize \textbf{$y\to 0$ καθώς $x\to-\infty$}};
        \draw[->, >=stealth, black, thick] (axis cs:-1.5, 0.4) -- (axis cs:-2.5, 0.1);
        
        % Point (1,1)
        \addplot[mark=*, mark options={fill=red, scale=1.2}] coordinates {(1,1)};
        \node[anchor=south east] at (axis cs:0.9, 1.1) {\footnotesize \textbf{$f(1)=1$}};
        
        \node[below left] at (axis cs:0,0) {$0$};
    \end{axis}
\end{tikzpicture}
\end{center}

Mε βάση το εμπλουτισμένο σχήμα:

\begin{erwthma}

\item Η γραφική παράσταση της $f$ αποτελείται από δύο διακριτούς κλάδους, στα αριστερά και στα δεξιά του $x = 2$, όπου και σχηματίζεται κατακόρυφη ασύμπτωτη. Συνεπώς, εφόσον η συνάρτηση δεν ορίζεται στο $x=2$, το πεδίο ορισμού είναι:
\[ \varDelta_f = \mathbb{R} \setminus \{2\} = (-\infty, 2) \cup (2, +\infty) \]
Από το γράφημα, διαβάζουμε ευθέως τις οριακές καταστάσεις για τα ζητούμενα όρια:
\begin{itemize}
    \item Καθώς το $x$ πλησιάζει το $2$ από αριστερά ($x \to 2^-$), η καμπύλη ανεβαίνει διαρκώς (προς τη θετική κατεύθυνση του άξονα $y'y$), άρα: $\lim\limits_{x \to 2^-} f(x) = +\infty$.
    \item Καθώς το $x$ πλησιάζει το $2$ από δεξιά ($x \to 2^+$), η δεξιά καμπύλη ομοίως ανεβαίνει προς το άπειρο, άρα: $\lim\limits_{x \to 2^+} f(x) = +\infty$.
    \item Καθώς το $x \to -\infty$, ο αριστερός κλάδος τείνει να ταυτιστεί με τον άξονα $x'x$ (οριζόντια ασύμπτωτη $y=0$), άρα: $\lim\limits_{x \to -\infty} f(x) = 0$.
    \item Καθώς το $x \to +\infty$, ο δεξιός κλάδος τείνει ασυμπτωτικά στην οριζόντια διακεκομμένη ευθεία (οριζόντια ασύμπτωτη $y=1$), άρα: $\lim\limits_{x \to +\infty} f(x) = 1$.
\end{itemize}

\item Ζητείται το όριο $\lim\limits_{x \to 2} e^{\frac{1}{f(x)}}$. 
Εφόσον δείξαμε στο πρώτο ερώτημα ότι και τα δύο πλευρικά όρια στο 2 ταυτίζονται στο $+\infty$, ισχύει συνολικά:
\[ \lim_{x \to 2} f(x) = +\infty \]
Οπότε, για τον εκθέτη προκύπτει:
\[ \lim_{x \to 2} \frac{1}{f(x)} = 0 \]
Η εκθετική συνάρτηση είναι συνεχής, άρα μπορούμε να θέσουμε το όριο του εκθέτη εντός, και υπολογίζουμε:
\[ \lim_{x \to 2} e^{\frac{1}{f(x)}} = e^0 = 1 \]

\item Από τον κάναβο στο γράφημα, παρατηρούμε καθαρά ότι η καμπύλη διέρχεται από το σημείο με συντεταγμένες $(1, 1)$. Οπότε $f(1) = 1$. 
Καθώς η συνάρτηση δίνεται ως συνεχής (στο πεδίο ορισμού της), ισχύει ότι το όριο της $f$ στο σημείο $1$ συμπίπτει με την τιμή της: $\lim\limits_{x \to 1} f(x) = f(1) = 1$.
Άρα:
\[ \lim_{x \to 1} \ln(f(x)) = \ln\left(\lim_{x \to 1} f(x)\right) = \ln(1) = 0 \]

\item Εξετάζοντας την πορεία της γραφικής παράστασης (κλίσεις), καταγράφουμε τα εξής στοιχεία μονοτονίας:
\begin{itemize}
    \item Για $x \in (-\infty, 2)$, η $f$ ξεκινά από το $0$ (καθώς $x \to-\infty$) και αυξάνεται σταθερά μέχρι να τείνει στο $+\infty$ (στο $2^-$). Συνεπώς, είναι \textbf{γνησίως αύξουσα} σε αυτό το διάστημα.
    \item Για $x \in (2, +\infty)$, η $f$ μειώνεται διαρκώς από το $+\infty$ (στο $2^+$) καταλήγοντας ασυμπτωτικά στο $y=1$. Άρα, είναι \textbf{γνησίως φθίνουσα} σε αυτό το διάστημα.
\end{itemize}
Εφόσον η μονοτονία δεν μεταβάλλεται \textit{εντός} των συνεχών διαστημάτων πεδίου ορισμού (δηλαδή δεν έχουμε καμία μετάβαση από αύξουσα σε φθίνουσα ή το αντίστροφο μέσα σε ένα ενιαίο διάστημα), η συνάρτηση \textbf{δεν παρουσιάζει κανένα (τοπικό ή ολικό) ακρότατο}. Η αλλαγή από αύξουσα σε φθίνουσα λαμβάνει χώρα εκατέρωθεν του $x=2$ (στο οποίο υπάρχει κατακόρυφη ασύμπτωτη και σημείο ασυνέχειας άρα και εξαιρείται), οπότε δεν ορίζεται ακρότατο εκεί.

\end{erwthma}
\end{thema}
%=========== ΤΕΛΟΣ ΛΥΣΗΣ - ΘΕΜΑ Β_89 ===========
